\subsection{Brakerski-Vaikuntanathan}
Brakerski and Vaikuntanathan~\cite{cryptoeprint:2011/344} first proposed decomposing ciphertext elements into small digits as a means of controlling the noise growth \textcolor{red}{of this operation}. The idea has since been generalized into the idea of \emph{gadget decomposition} and applied in a variety of forms, including digit and CRT (RNS) decomposition. Following the convention of~\cite{cryptoeprint:2021/204, cryptoeprint:2022/915}, we refer to this family of decomposition-based techniques collectively as \emph{BV}.

\begin{definition}[Digit Decomposition]
    Given an integer $\gamma\in\mathbb{Z}_q$ for a modulus $q\geq 2$, base $\beta$ and decomposition parameter $\ell$, the digit decomposition $\decomp{\gamma}{} = (\gamma_0, \gamma_1, \dots, \gamma_{\ell-1}) \in \mathbb{Z}_\beta^\ell$ outputs an ordered list of $\ell$ base-$\beta$ digits satisfying
    \begin{equation*}
        \gamma = \gamma_0\frac{q}{\beta} + \gamma_1\frac{q}{\beta^2} + ... \gamma_{\ell - 1}\frac{q}{\beta^{\ell}} = \sum_{i=0}^{\ell-1}\gamma_i\frac{q}{\beta^{i+1}}
    \end{equation*} 
\end{definition}

%%
%% Why its helpful
Digit decomposition has two important properties. First, if the digits $\gamma_i$ of $\decomp{\gamma}{}$ are small, then $\langle \decomp{\gamma}{}, \epsilon\rangle \leq \ell\cdot \max\{\gamma_i\cdot \epsilon_i\} \leq \beta \epsilon \ll \gamma \epsilon$. This is the property that allows some control over the noise growth. Second, two integers can be multiplied through the decomposition. By construction, $\langle\decomp{\gamma}{}, g\rangle = \gamma$ for the fixed vector $g = (q/\beta, \dots, q/\beta^{\ell})$, meaning $\langle\decomp{\gamma}{}, g\cdot x\rangle = \gamma\cdot x$ for an arbitrary integer $x\in \mathbb{Z}$. For either of these properties, the particular form of $g$ is unimportant as long as the digits $\gamma_i$ remain small. Furthermore, as we are generally working with some intrinsic error in FHE, \cite{cryptoeprint:2018/421} proposed allowing some error $\epsilon$ in the reconstruction, giving the definition from their paper verbatim:

%%
%% Gadget Decomposition
\begin{definition}[(Abstract) Gadget Decomposition]
    \label{def:tfhe:gadget-decomposition}
    An efficient algorithm $\decomp{v}{g,\beta,\epsilon}$ is a valid decomposition on the gadget $g\in \mathbb{Z}_q^\ell$ with quality $\beta \in \mathbb{Z}^{+}$ and precision $\epsilon \in \mathbb{Z}^+$ if and only if for any \gls{rlwe} sample $v\in R_Q$ it efficiently and publicly outputs a small vector $u \in \mathbb{Z}_q$ such that $\|u\|_\infty \leq \beta$ and $\langle u, g\rangle \leq v + \epsilon$.

    % Furthermore, the expectation of u · H − v must to be equal to 0 when v is uniformly distributed in M .
\end{definition}

\begin{itemize}
    \item Other (equivalent) form more convenient to work with, where the output of $\decomp{v}{g,\beta,\epsilon} \to D(v)$ and $g\cdot u \to P(u)$ are considered.
\end{itemize}

\begin{definition}[Gadget Property]
    Two vectors $P(x) \in \mathbb{Z}_Q^\ell$ and $D(y) \in \mathbb{Z}_\beta^\ell$ are said to have the \textit{gadget property} if and only if they satisfy \eqref{def:gadget-property} for any pair of \gls{rlwe} samples $(x,y)\in R_Q^2$ with precision $\epsilon \in \mathbb{Z}_q^+$ and quality $\beta$ where $\|D(y)\|_\infty \leq \beta$. If $\epsilon = 0$ then $P(x)$ and $D(y)$ are said to have the \textit{exact} gadget property.
    \begin{equation}
        \label{def:gadget-property}
        \langle P(x), D(y)\rangle = x\cdot y + \epsilon \pmod{Q}
    \end{equation}
\end{definition}

% \begin{itemize}
%     \item Can be applied component-wise to polynomials $\gamma\in R_q$ instead of $\mathbb{Z}_q$
%     \item Centered representation
% \end{itemize}

\begin{lemma}%[Layered Gadget Decomposition]
    \label{thm:gadget-layers}
    The composition of two gadget decompositions is itself a valid gadget decomposition.
    \begin{gather*}
        \langle P_1(u), D_1(v)\rangle = u\cdot v + \epsilon_1, \quad \langle P_2(u), D_2(v) \rangle = u\cdot v + \epsilon_2 \pmod{Q} \\
        P_1\circ P_2 = P_1(P_2(u)) \implies \langle P_1\circ P_2, D_1 \circ D_2 \rangle = u\cdot v + 2\epsilon \pmod{Q} 
    \end{gather*}

    \begin{proof}
        This is clear from applying the gadget property twice.
        \begin{gather*}
            \langle P_1(P_2(u)), D_1(D_2(u)) \rangle = P_2(u) \cdot D_2(u) + \epsilon_1 = u\cdot v + \epsilon_2 + \epsilon_1 \pmod{Q}
        \end{gather*}
    \end{proof}
\end{lemma}






\subsubsection{RNS Instantiation}
Digit decomposition cannot be applied directly the \gls{rlwe} sample in an \gls{rns} scheme, where the system only has access to the residues. When Bajard et. al. first implemented the BV cryptosystem using \gls{rns} \cite{cryptoeprint:2016/510} they exchanged digit decomposition for \gls{crt} decomposition, exploiting the fact that this decomposition is already required by the \gls{rns}.

%% TODO: Clean up, make sure \hat{q} notation is known 
By stating the \gls{crt} as an inner product \autoref{eq:crt-gadget}, it is immediately clear that \gls{crt} with moduli $q_i$ is an exact decomposition on $g = (\hat{q}_0\cdot\hat{q}_0^{-1}\bmod{Q}, \dots, \hat{q}_{k-1}\cdot\hat{q}_{k-1}^{-1}\bmod{Q})$ with $\beta = \max_i\|q_i\|_\infty/2$ (assuming centered residues) and $\epsilon = 0$. 

\begin{itemize}
    \item CRT idempotent used in \cite{cryptoeprint:2016/510}
    \item HPS optimization for specifically keyswitching \cite{cryptoeprint:2018/117}
    \item Definition: gadget with $\beta = Q/2, \epsilon = 0$
\end{itemize}

\begin{align}
    \hps{P}(u) = () \\
    \hps{D}(v) = ()
\end{align}

\begin{itemize}
    \item Note on HPS gadget = RNS limb lookup $\to$ \autoref{sec:impl:core}
    \item HPS is exact but noise is large ($\beta = Q/2,\; \epsilon = 0$); can apply digit decomposition to each residue polynomial as per \autoref{thm:gadget-layers}, lowering $\beta \to B/2$ and still exact ($\epsilon = 0$) if $B = \lfloor \max{q_i}/\ell \rfloor + 1$
\end{itemize}